\subsection*{前言}
\addcontentsline{toc}{subsection}{前言}

本书整理自 UC Berkeley / Berkeley RDI 课程《CS294/194-280: Advanced Large Language Model Agents, Spring 2025》的公开视频、官方 slides、课程页与 supplemental readings，并在最终定稿时进一步吸收了 Berkeley Fall 2024《Large Language Model Agents》、Berkeley Fall 2025《Agentic AI》以及 Stanford CS329A《Self-Improving AI Agents》的官方课程材料。它讨论的不是单一模型技巧，而是一整套关于高级 LLM agents 的问题框架：推理时计算（inference-time computation）应如何分配，后训练（post-training）怎样改变 agent 的判断质量，记忆与规划如何支持长程决策，工具与环境反馈怎样把语言模型推进到 coding、GUI、web、robotics 和 formal mathematics，最后又该如何在更强能力出现时建立安全与安全边界。

这本书刻意不采用“看完视频后写一份摘要”的做法。每一讲在进入全书之前，都先作为独立 lecture workspace 构建 source manifest、transcript ledger、slide parsing、coverage units、figure provenance、reading integration、evaluator report 与 validator report；只有通过 lecture gate 的章节，才允许合并进入 textbook。换句话说，本书首先是一个有审计轨迹的教材工程（harness-managed textbook build），其次才是一份成书后的阅读对象。

这种做法有两个直接后果。第一，本书比普通课程笔记更重视覆盖完整性，而不是篇幅压缩。重要定义、公式、算法流程、代码逻辑、实验设置、反例、失败模式与边界条件，原则上都应先被记录，再考虑如何组织表达。第二，本书明确区分三类内容：来自课程源的事实，基于多源材料的整合解释，以及为帮助读者理解而做的延伸性说明。前两类必须能回查 provenance；第三类若出现，则应服从原始材料而不是替代原始材料。

从课程结构上看，这门课并不是把“agent”当作一个宽泛标签来使用。前半部分从推理、判断学习、记忆与规划出发，逐步建立 agent 的能力结构；中段进入 post-training recipe、coding agent 与 multimodal environments，把抽象能力放进具体工具链和交互环境中考察；后半部分则转向 formal mathematics、autoformalization、theorem proving、abstraction/discovery，以及 safe and secure agentic AI。把这些讲次合并成书之后，读者应看到的是一条连续主线：高级 agent 的本质在于如何把模型、搜索、反馈、工具、环境与边界条件组织成一个可执行系统。扩展部分再把这条主线向 Berkeley 2024/2025 与 Stanford CS329A 拉开，补齐系统工程、agent evaluation、自改进和最新课程生态中的公开材料。

本书的目标读者不是只想快速浏览结论的人，而是希望在不依赖原视频的情况下系统掌握该课程的研究问题、技术机制和方法边界的读者。你可以把它当作自学教材、课程复习材料、研究综述入口，或者作为定位具体 lecture provenance 的索引。若某一章的个别细节因为字幕质量、slides 信息密度或 reading 可得性而仍存在不确定性，本书不会掩盖这一点，而会把低置信度片段、遗漏原因和未解决问题保留在对应的 sidecar 与附录中。

因此，这本书想提供的不是“替你看完视频”的轻量总结，而是一份可读、可查、可复核、可继续增量更新的课程教材基座。
